
Once a group has been created as in 
Sections~\ref{sec:generic:permutation:groups} and~\ref{sec:lineargroups}, 
a group theoretic activity can be performed. 
For this purpose, Orbiter provides the \verb'-group_theoretic_activity' option. 
Tables~\ref{tab:grouptasks}-\ref{tab:grouptasks5} list the available commands.
\orbitertableofcommands{group_theoretic_activity_1}
\orbitertableofcommands{group_theoretic_activity_2}
\orbitertableofcommands{group_theoretic_activity_3}
\orbitertableofcommands{group_theoretic_activity_4}
\orbitertableofcommands{group_theoretic_activity_5}
Table~\ref{tab:grouptasks:elementprocessing} lists commands for element processing. 
\orbitertableofcommands{element_processing}


\bigskip

The command

\sourcecode{PGL_3_2_order_invariant}

computes the order invariant of $\PGL(3,2).$ 
This is the statistic of the orders of the elements in the group. 
The following latex output will be created:

\orbiteroutput{GRAPHICS/LATEX/PGL_3_2_order_invariant}


This shows us that the group contains 48 elements of order 7, 
42 elements of order 4 and so forth.
The order invariant is invariant under group isomorphisms, 
and hence can be used to distinguish non-isomorphic groups. 
If the order invariants of two groups differ, the groups are not isomorphic.

\bigskip


The command

\sourcecode{PGL_3_2_elements}

creates all elements of $\PGL(3,2)$ 
and writes them into the file \verb'PGL_3_2_elements.csv'.

\bigskip


The command

\sourcecode{Sym_3_elements}

creates a report of the elements of the group $\Sym(3)$: 
\orbiteroutput{GRAPHICS/LATEX/report_group_elements_Sym3}


\bigskip


The command

\sourcecode{Sym_3_make_element_tree}

produces a tree of all group elements of $\Sym(3)$ according to the selected stabilizer chain. 
\orbiteroutput{GRAPHICS/WRAPPERS/sym_3_elements_tree}
The stabilizer chain is defined by a sequence of points in the permutation domain called a base.
For $\Sym(n)$ the default base is the sequence $0,1,2,\ldots,n-1.$


\bigskip



It is possible to compute all powers of a fixed element, as in the following command:

\sourcecode{Cycle_12_power}

We create the 12 powers of the cycle 
$$
(0,1,2,3,4,5,6,7,8,9,10,11).
$$
The output is
\orbiteroutput{GRAPHICS/LATEX/report_group_element_powers}


\bigskip


The command

\sourcecode{PGL_3_4_singer}

finds all Singer cycles in $\PGL(3,4)$ 
whose matrix is the companion matrix of a polynomial. 
The first one found is 


\orbiteroutput{GRAPHICS/LATEX/report_group_singer}


whose projective order is $21.$ 
Here, we are using the numeric form of field elements, 
so $2$ is $\omega$ and $3$ is $\omega+1.$

\bigskip

Suppose we want to multiply two group elements. 
The following command shows an example in $\GL(2,8).$ 
The elements of $\GL(2,8)$ 
are coded as vectors of length 4 with entries coded as integers in the interval $[0,7]$. 
We multiply the elements coded by 0,1,2,3 and 4,5,6,7:

\sourcecode{GL_2_8_multiply}

The output is
\orbiteroutput{GRAPHICS/LATEX/report_group_multiply_GL_2_8}
Note that the output shows the codings of the three group elements. 
This way, the result of this computation can be processed further easily.
The same example over $\bbF_7$ is:

\sourcecode{GL_2_7_multiply}

The output is
\orbiteroutput{GRAPHICS/LATEX/report_group_multiply_GL_2_7}
We can compute the inverse of a group element:

\sourcecode{GL_2_7_inv}

The output is

\orbiteroutput{GRAPHICS/LATEX/report_group_inverse_GL_2_7}

We can raise a group element to a power:

\sourcecode{GL_2_7_power}

The output is

\orbiteroutput{GRAPHICS/LATEX/report_group_power_GL_2_7}

\bigskip



Suppose we want to find a small generating set for a group. 
The following command invokes a randomized algorithm to find a generating set of 2 elements for the 
group $\PGGL(4,4).$ The command specfies the size of the desired generating system and an upper bound for the number of attempts. 


\sourcecode{PGGL_4_4_small_generating_set}


\bigskip




The next example computes the action of 
a specific group element on the set of planes through a line. 
The planes have been computed in Section~\ref{sec:grassmannian}.

\sourcecode{on_planes}

The output is

\orbiteroutput{GRAPHICS/LATEX/report_group_on_planes}

\bigskip

The command \verb'-random_element' 
can be used to pick a random element in a group.
Here is an example. 
We create the group 
$\PGL(4,2)$ and then ask for a random element:

\sourcecode{PGL_4_2_random_element}


The output is the following csv file (the actual content may differ, as the output is pseudo-random).

\orbiteroutput{GRAPHICS/WRAPPERS/random_element}

The contents is as follows:
\begin{enumerate}
\item
The column \verb'Rk' is the rank number of the group element in the group. It is a number between zero and group order minus one. 
\item
The column \verb'Path' is the coset decomposition in the basic orbits determined by the rank. 
\item
The column \verb'Tl' is the vector whose $i$-th entry is the length of the $i$-th basic orbit. 
\item
The column \verb'ElementLatex' is a latex description of the group element.
\item
The column \verb'ElementCoded' is the coded description of the group element.
\item
The column \verb'Perm' is the permutation action of the group element in the natural action.
\item
The column \verb'CycleType' is the cycle type of the permutation action of the group element in the natural action.
\item
The column \verb'Sign' is the sign of the permutation action of the group element in the natural action (even=1, odd=-1).
\item
The column \verb'Order' is the order of the group element in the natural action.
\end{enumerate}




\bigskip




The command

\sourcecode{PGL_2_3_subgroup_lattice}

computes the subgroup lattice of 
$\PGL(2,3)$ and writes the data to a csv file called 
$$
\verb'PGL_2_3_subgroup_lattice.csv'.
$$

The command

\sourcecode{PGL_2_3_subgroup_lattice_load_and_draw_by_groups}

reads this file and produces a drawing of the lattice by groups. 
The nodes correspond to subgroups. 
The lattice is ``upside down,'' 
i.e. the unit subgroup is at the top and the full group is at the bottom. 
The maximal inclusions are drawn as edges between the nodes:

\orbiteroutput{GRAPHICS/WRAPPERS/sgl_PGL_2_3_by_groups}


\bigskip


The command

\sourcecode{PGL_2_3_subgroup_lattice_load_and_draw_by_orbits}

also reads the csv file and produces a drawing of the lattice by orbits. 
The nodes correspond to conjugacy classes of subgroups. 
The nodes are labeled by symbols of the form  
$$
{a \atop b}.
$$
Here, $a$ is the order of the subgroups 
in the orbit and $b$ is the number of conjugate groups.  

\orbiteroutput{GRAPHICS/WRAPPERS/sgl_PGL_2_3_by_orbits}


\bigskip

The command

\sourcecode{PGL_3_13_2B_RNF}

computes the rational normal form of the matrix
$$
A=\left[
\begin{array}{ccc}
1 & 1 & 11\\
1 & 6 & 12\\
7 & 10 & 5\\
\end{array}
\right]
$$
in $\PGL(3,13)$.
The rational normal form is 
$$
R = \left[
\begin{array}{ccc}
12 & 0 & 0\\
0 & 12 & 0\\
0 & 0 & 1\\
\end{array}
\right].
$$
A matrix $B$ such that $B^{-1}AB=R$ is 
$$
B = \left[
\begin{array}{ccc}
7 & 12 & 9\\
12 & 0 & 11\\
0 & 12 & 12\\
\end{array}
\right].
$$

\bigskip


